\section{Explicit reconstruction of the inherited section multipliers}
\label{sec:sm-section}

The occupation observable is the indicator of the actual section
$Y_R^*$. The historical native section-multiplier lemma and module
84 already prove its strong multiplier property, entire occupation
phase and weak parameter continuity; modules 84--86 contain the
inherited full fixed-band spectral and arithmetic arguments.
Those results are retained, with their stated geometric and analytic
hypotheses. This section supplies a fully explicit reconstruction
of the section interface and its external hypotheses. It is a
proof-completeness addition, not a claim that those multiplier or
fixed-band theorems are new. No section side is rounded and no
return event is changed.

Use the common coordinates $(\alpha,\varphi)$, with
$p=\sin\varphi$, and the radius-dependent arclength
$s=R\alpha$. The classical stable cones in arclength coordinates
have slopes between $-K$ and $-a$, where the uniform choices
\begin{equation}\label{eq:sm-cone-constants}
 a=100/47,\qquad K=170/9
\end{equation}
follow from the actual curvature and flight bounds. In the common
coordinates every such cone is contained in
\begin{equation}\label{eq:sm-common-cone}
 -\overline K\le d\varphi/d\alpha\le-\underline a,
 \qquad \underline a=45/47,\qquad \overline K=799/90.
\end{equation}
The common cone contains the physical stable directions. The test
arcs used below are tangent to this cone, have the fixed curvature
and short-length bounds, and are closed under restriction. Invariance
of the enlarged test family under every $T_R$ is a separate spectral
requirement. The sign is reversed from the unstable curves in
Section~\ref{sec:lc-length}. The change $s=R\alpha$ and its inverse
have uniformly bounded derivatives on the given radius interval.

We use precisely the stable-curve norm construction of
Demers--Zhang~\cite[Sect.~3.3, Lemmas~5.2--5.3]{DZUniform}.
In that construction fix common admissible exponents
$\alpha_{\rm DZ},\beta_{\rm DZ},p_{\rm DZ},q_{\rm DZ}$ and
curve length and curvature bounds, with
$0<q_{\rm DZ}<p_{\rm DZ}\le1/3$,
$0<\alpha_{\rm DZ}<1-\varsigma_0$ for the fixed summability
parameter $\varsigma_0\in(0,1)$ of that construction, and
$0<\beta_{\rm DZ}\le\min\{\alpha_{\rm DZ},p_{\rm DZ}-q_{\rm DZ}\}$.
The norms are denoted $\|\cdot\|_{\mathcal B}$ and
$|\cdot|_{\mathcal B_w}$ on the common cylinder. The strong norm
dominates the stable norm and the weak norm tests homogeneous stable
arcs against functions with $C^{p_{\rm DZ}}$ norm at most one.
Here $C^p$ has the cited little H\"older convention, the closure
of smooth functions in the H\"older norm on each test arc.
The stable part has the following defining consequence: on an
admissible subarc $V$,
\begin{equation}\label{eq:sm-short-arc}
 \left|\int_V h\psi\,dm_V\right|
 \le \|h\|_{\mathcal B}|V|^{\alpha_{\rm DZ}}
                           |\psi|_{C^{q_{\rm DZ}}(V)}.
\end{equation}
The strong multiplier proof also applies in each radius-specific
Demers--Zhang arclength space, with common norm parameters and curve
bounds. Its constants include the uniformly bounded coordinate
scalings. The weak difference of two moving sections is asserted
on the fixed common test-curve construction just specified; no
identification of different dynamical spaces is implicit. In the
measure embedding of this construction a function $h$ represents
$h\nu$; the constant function $1$ therefore represents $\nu$.
The coordinate expressions are
$d\nu=\cos\varphi\,d\alpha\,d\varphi/(4\pi)
=\cos\varphi\,ds\,d\varphi/(4\pi R)$.
These expressions describe the reference measure and are not a
replacement of $h$ by a density with respect to flat area. In
particular no division by $\cos\varphi$ is used at grazing.
Lemma~3.4 and the measure embedding of the cited construction
make the mass functional $\ell(h)=\int h\,d\nu$ continuous on
$\mathcal B_w$, by testing against the constant function $1$.
The assertions concern these norm constructions. They
do not assert that the complete jointly twisted collision operators
already satisfy the separate hypotheses (H1)--(H5), (C1)--(C4), or
the spectral conclusions of that reference.

\begin{lemma}[Finite transverse partition for the actual section]
\label{lem:sm-grid}
For every $R\in[R_-,R_+]$ there is a modulo-zero partition of the
collision cylinder into at most $121$ open simply connected cells
on each of which $\eta_R=\one_{Y_R^*}$ is constant. Each cell has
at most four smooth sides. Every side is uniformly transverse to
\eqref{eq:sm-common-cone}; a lower bound for the sine of the angle is
\begin{equation}\label{eq:sm-transversality}
 \vartheta_{\rm sec}=
 \frac{\underline a}{\sqrt{1+\overline K^2}}>0.
\end{equation}
Every homogeneity strip meets at most $121$ cells. The original
section is contained in $|p|<5/6$, and its four moving vertical
sides have speed at most $6$ in $\alpha$, while its four moving
horizontal sides have speed less than $11$ in $\varphi$.
\end{lemma}
\begin{proof}
The four squares and the rectangle $U$ have at most ten distinct
$\alpha$ side values and ten distinct $p$ side values. Cut the
angular circle at the common seam $\alpha=\pi$, which lies outside
all these rectangles for the entire radius range, and extend the
coordinate sides into full coordinate lines. At most eleven angular
intervals and eleven $p$ intervals result. Removing the lines gives
the asserted partition. A cell remains a rectangle after
$p=\sin\varphi$, so it is simply connected and has at most four
sides. Empty or coincident cells reduce the count.

A unit tangent to a stable arc has components proportional to
$(1,v)$ with $-\overline K\le v\le-\underline a$. Its angle with
a vertical line has sine at least $(1+\overline K^2)^{-1/2}$;
its angle with a horizontal line has sine at least
$\underline a(1+\overline K^2)^{-1/2}$. This proves
\eqref{eq:sm-transversality}, also at the artificial seam. The
homogeneity-strip count is bounded by the total cell count.

Here is an explicit check of the moving sides. Write $F(R,q)$ for
the polynomial with coefficient $\sqrt3$ defining $q_R$ in the
physical model. For $49/100\le q\le53/100$,
\[
 F_q=8R(1+3q^2)-4+12q^2+\sqrt3(-12q+4q^3)<-2.
\]
Indeed its first three terms are at most
$8(47/100)(1+3(53/100)^2)-4+12(53/100)^2$, and its last term
is at most $(17/10)(-12(49/100)+4(53/100)^3)$: the expression
in parentheses is negative and $\sqrt3>17/10$. Their sum is
$-2.6842844$. Also $0<F_R\le8(53/100)(1+(53/100)^2)<6$.
Implicit differentiation therefore gives $0<q'_R<3$ and
$|\theta'_R|=2q'_R/(1+q_R^2)<6$. These estimates hold on the
whole given root interval, including its closed endpoints used for
the bounds; uniqueness of the physical root is already part of
the retained geometry.

The two moving squares have angular center speed $|\theta'_R|$
and momentum center speed $|(\sin\theta_R)'|$, both at most $6$.
Furthermore
\[
 \sin\theta_R+h_0\le
 \frac{2(53/100)}{1+(53/100)^2}+\frac1{200}<\frac56.
\]
All the other section momentum coordinates satisfy this same cap.
Thus every moving horizontal side in $\varphi$ coordinates has
speed at most $6/\sqrt{1-(5/6)^2}=36/\sqrt{11}<11$.
The initial normal square, the square centered at $(\pi/6,-1/2)$,
and $U$ are fixed. There are exactly four possibly moving sides
of each orientation. No assertion of independence of their motion
or of orbit visits has entered the calculation.
\end{proof}

\begin{theorem}[Strong multiplication by the exact section and occupation phase]
\label{thm:sm-multipliers}
Let $\mathcal B$ have the Demers--Zhang norm construction just
specified. There is a constant $C_{\rm sec}$ independent of $R$
such that multiplication by $\eta_R$ and by
$f_{R,z}=\exp(z\eta_R)$ is bounded on $\mathcal B$, with
\begin{equation}\label{eq:sm-strong}
 \|M_{\eta_R}\|_{\mathcal B\to\mathcal B}\le C_{\rm sec},
 \qquad
 \|M_{f_{R,z}}\|_{\mathcal B\to\mathcal B}
       \le C_{\rm sec}\max\{1,e^{\operatorname{Re}z}\}.
\end{equation}
For every real occupation frequency $\zeta$ the second bound is
$C_{\rm sec}$ at $z=i\zeta$. The map $z\mapsto M_{f_{R,z}}$
is entire in operator norm, and, for every integer $r\ge1$,
\begin{equation}\label{eq:sm-analytic}
 M_{f_{R,z}}=I+(e^z-1)M_{\eta_R},\qquad
 \partial_z^r M_{f_{R,z}}=e^zM_{\eta_R}.
\end{equation}
The constant depends on the fixed norm parameters, the common
curve bounds, and the explicit transverse partition constants,
but on no collision count or observation window.
\end{theorem}
\begin{proof}
Choose the H\"older exponent used in Lemma~5.3 of the cited
author version as
$\gamma=\max\{p_{\rm DZ},2\beta_{\rm DZ}+\epsilon_{\rm DZ}\}$
with one admissible $\epsilon_{\rm DZ}>0$. On every cell of
Lemma~\ref{lem:sm-grid}, both multipliers are constant. Their
piecewise $C^\gamma$ norms are therefore at most $1$ and
$\max\{1,e^{\operatorname{Re}z}\}$, respectively.
The partition is modulo zero, open and simply connected; its
boundary count, transversality and cells per homogeneity strip
have all been checked explicitly. Lemmas~5.2--5.3 consequently
give the two bounded multiplication operators on the strong space.
All parameters entering their estimates are uniform here. In
particular no factor $|\zeta|$ is incurred for the piecewise constant
occupation phase.

On smooth densities the identity
$e^{z\eta_R}=1+(e^z-1)\eta_R$ holds almost everywhere. The bounded
extensions and density of the original smooth class in
$\mathcal B$ prove the first operator identity in
\eqref{eq:sm-analytic}. Its norm-convergent power series proves
the asserted entire dependence and every derivative formula.
This argument invokes the precise multiplier theorem, not its
spectral corollaries.
\end{proof}

\begin{theorem}[Weak perturbation estimate for the moving physical section]
\label{thm:sm-moving}
Use the common stable-curve realization above. There is a constant
$C_{\rm move}$, independent of $R,S\in[R_-,R_+]$, such that
\begin{align}
 \|M_{\eta_R}-M_{\eta_S}\|_{\mathcal B\to\mathcal B_w}
       &\le C_{\rm move}|R-S|^{\alpha_{\rm DZ}},
                                                \label{eq:sm-weak-move}\\
 \|M_{f_{R,z}}-M_{f_{S,z}}\|_{\mathcal B\to\mathcal B_w}
       &\le |e^z-1|C_{\rm move}|R-S|^{\alpha_{\rm DZ}}.
                                                \label{eq:sm-phase-move}
\end{align}
The measures and section normalizing constant are the original
ones. These inequalities concern multiplication only; a moving
collision map has its own additional perturbation estimates.
\end{theorem}
\begin{proof}
Write $d=|R-S|$. For $d=0$ both assertions are identities. The
symmetric difference of the two sections is contained in the union
of strips swept out by the eight moving sides. Their coordinate
widths are at most $11d$, by Lemma~\ref{lem:sm-grid}. Along any
stable arc satisfying \eqref{eq:sm-common-cone}, both coordinates
are strictly monotone in arclength and their absolute speeds are
at least $\vartheta_{\rm sec}$. The intersection with any one
of these strips has length at most $11d/\vartheta_{\rm sec}$.

Cut an admissible short stable arc at the coordinate sides of both
section grids and, if necessary, at the angular seam. There are
at most twenty side values of each orientation. On each chart
piece a monotone arc therefore has at most forty-one complementary
intervals; allowing a seam split gives the safe bound eighty-two.
The difference $\eta_R-\eta_S$ is constant on each such interval,
with value $0$, $1$, or $-1$. Every nonzero interval is contained
in a swept strip. One way to see the last assertion is to refine
also by the strip endpoints, which are among the same two sets
of side values: membership in each strip is then constant on a
component. Thus all nonzero subarcs have length at most
$11d/\vartheta_{\rm sec}$. Curvature, cone membership and
homogeneity are preserved by restriction.

For a weak test $\psi$ of $C^{p_{\rm DZ}}$ norm at most one,
its restricted $C^{q_{\rm DZ}}$ norm is bounded by a uniform
constant $C_{pq}$, since $p_{\rm DZ}\ge q_{\rm DZ}$ and the
maximum curve length is fixed. Applying
\eqref{eq:sm-short-arc} on each nonzero interval gives
\[
 \left|\int_W h(\eta_R-\eta_S)\psi\,dm_W\right|
 \le82 C_{pq}\left(\frac{11}{\vartheta_{\rm sec}}\right)^{\alpha_{\rm DZ}}
             d^{\alpha_{\rm DZ}}\|h\|_{\mathcal B}.
\]
This is first proved for smooth densities, for which restrictions
and pointwise products are literal. Taking the weak norm and then
using the bounded strong multipliers of
Theorem~\ref{thm:sm-multipliers} extends it to their strong-space
completion. The continuous embedding into the weak space makes
the extension the same operator. This proves
\eqref{eq:sm-weak-move}; the first identity in
\eqref{eq:sm-analytic} proves \eqref{eq:sm-phase-move}.
\end{proof}

\begin{corollary}[Exact occupation and endpoint tests in a qualified transfer formula]
\label{cor:sm-transfer}
Suppose, separately, that the physical operator
\[
 \mathcal A_{R;\xi,s}h=\mathcal L_R
       \bigl(e^{i(\xi\cdot\kappa_R+s\tau_R)}h\bigr)
\]
has a bounded realization on this strong space at the displayed
parameters. Define
$\mathcal L_{R;\xi,\zeta,s}
 =\mathcal A_{R;\xi,s}M_{f_{R,i\zeta}}$.
For every finite collision count $m$ for which the realization is
available, its physical moment is exactly
\begin{equation}\label{eq:sm-exact-transfer}
 \frac1{c_*}\ell\left(
 M_{\eta_R}\mathcal L_{R;\xi,\zeta,s}^{m}
                         (\eta_R\nu)\right)
 =\int \eta_R(T_R^mx)
 e^{i(\xi\cdot K^{\rm c}_{m,R}+\zeta A_{m,R}+sL_{m,R})}
                 \,d\nu_R^*(x).
\end{equation}
The section multipliers in this formula satisfy
\eqref{eq:sm-strong}, and the dependence on $\zeta$ is entire.
\end{corollary}
\begin{proof}
The physical transfer convention is
$(\mathcal L_Rh)(b)=h(b\circ T_R)$ on measures.
An induction shows that $m$ iterates multiply the trajectory by
the displayed flight and displacement phases and by
$\prod_{j=0}^{m-1}e^{i\zeta\eta_R(T_R^jx)}$.
The latter product is exactly $e^{i\zeta A_{m,R}}$.
The terminal multiplier tests $\eta_R(T_R^mx)$ and contributes
no extra occupation count. The initial measure is $\eta_R\nu$;
the single division by $c_*$ gives the original entrance law.
This proves the equality on the actual physical trajectories.
When the bounded realization is used, the multiplier theorem
justifies its section factors. The inherited fixed-band realization
and power estimates have their own proof bodies and hypotheses;
this multiplication argument supplies no extension of their
constants to a growing roof-frequency band.
\end{proof}

\paragraph{Reconstructed interface and retained spectral scope.}
The original section indicator, full occupation Fourier phase and
moving section cuts have explicit strong multiplier and weak
perturbation estimates with checked geometric hypotheses. This
reconstruction does not reopen the existing fixed-band author
proofs or replace them by a new assumption. Control on growing roof-frequency bands is a separate possible
spectral estimate and is not supplied by this multiplier calculation.
The retained fixed-window theorem will evaluate the scalar bridge
in Section~\ref{sec:ui-exact} without such an additional estimate.
The original endpoint shortage \eqref{eq:lc-final-shortage}, the
normalized local derivative charge on that alternative route, and
the actual microscopic denominator for conditional conclusions
remain at their stated scopes.
